An extension remembers both its constituents and the object they form.
The extension correspondence, with the quotient written first, is
\[
 \M_{\alpha}\times\M_{\beta}
 \xleftarrow{\ p_{\alpha,\beta}\ }\M_{\alpha,\beta}^{\mathrm{ext}}
 \xrightarrow{\ q_{\alpha,\beta}\ }\M_{\alpha+\beta},
 \qquad 0\longrightarrow E_\beta\longrightarrow E
 \longrightarrow E_\alpha\longrightarrow0.
\]
Coefficients describe the states carried by an object. Pullback puts two
incoming states on the extension space. Proper pushforward then produces a
complex on the space of middle objects; a further morphism must identify its
contribution to the chosen outgoing states.

Fix rational constructible complexes $K_\gamma$ on $\M_\gamma$, in the
analytic topology. For quotient stacks, use equivariant constructible
complexes, so that stabilizer actions are retained. Assume the displayed
pullbacks, external products and proper pushforwards preserve the chosen
derived category. All degree shifts are included in $K_\gamma$ and the
following degree-zero morphisms:
\[
 c_{\alpha,\beta}:Rq_{\alpha,\beta*}p_{\alpha,\beta}^*
 (K_\alpha\boxtimes K_\beta)\longrightarrow K_{\alpha+\beta}.
\]
No identification of these two complexes is part of the notation. For
$a\in\mathbb H^r(\M_\alpha,K_\alpha)$ and
$b\in\mathbb H^s(\M_\beta,K_\beta)$, define
\[
 a\star b=\mathbb H(c_{\alpha,\beta})\bigl(p_{\alpha,\beta}^*
 (a\boxtimes b)\bigr)
 \in\mathbb H^{r+s}(\M_{\alpha+\beta},K_{\alpha+\beta}).
\]
Here the class after pullback is viewed in the hypercohomology of
$Rq_{\alpha,\beta*}p_{\alpha,\beta}^*(K_\alpha\boxtimes K_\beta)$,
using the defining equality of derived global sections under pushforward.

The role of $c$ is already visible when every space is a point. Take
$K=\mathbb Qe\oplus\mathbb Qf$ in degree zero and set
\[
 c(e,e)=f,\qquad c(f,e)=e,\qquad c(e,f)=c(f,f)=0.
\]
All geometric pullbacks, pushforwards and base-change maps are identities.
Nevertheless $(e\star e)\star e=e$ and
$e\star(e\star e)=0$. Thus geometric composition alone leaves a real
condition on the coefficient morphism.

For three charges, let $\mathcal F_{\alpha,\beta,\gamma}$ be the stack of
filtrations $0\subset E_1\subset E_2\subset E_3$ with successive
constituents of charges $\gamma,\beta,\alpha$. Write
\[
 s:\mathcal F_{\alpha,\beta,\gamma}\longrightarrow
 \M_\alpha\times\M_\beta\times\M_\gamma,
 \qquad r:\mathcal F_{\alpha,\beta,\gamma}\longrightarrow
 \M_{\alpha+\beta+\gamma}
\]
for the maps remembering the ordered constituents and $E_3$. Assume every $q_{\alpha,\beta}$ and $r$ is proper, the forgetful squares
are Cartesian, and proper base change and external-product exchange give
the following identifications:
\[
 \begin{aligned}
 D&=Rr_*s^*(K_\alpha\boxtimes K_\beta\boxtimes K_\gamma),\\
 C_L:D&\xrightarrow{\sim}Rq_{\alpha+\beta,\gamma*}
 p_{\alpha+\beta,\gamma}^*
 \bigl((Rq_{\alpha,\beta*}p_{\alpha,\beta}^*
 (K_\alpha\boxtimes K_\beta))\boxtimes K_\gamma\bigr),\\
 C_R:D&\xrightarrow{\sim}Rq_{\alpha,\beta+\gamma*}
 p_{\alpha,\beta+\gamma}^*
 \bigl(K_\alpha\boxtimes(Rq_{\beta,\gamma*}p_{\beta,\gamma}^*
 (K_\beta\boxtimes K_\gamma))\bigr).
 \end{aligned}
\]
The order of the three factors is unchanged. Define two maps from $D$ to
$K_{\alpha+\beta+\gamma}$ by
\[
 \begin{aligned}
 L&=c_{\alpha+\beta,\gamma}\circ
 Rq_{\alpha+\beta,\gamma*}p_{\alpha+\beta,\gamma}^*
 (c_{\alpha,\beta}\boxtimes1)\circ C_L,\\
 R&=c_{\alpha,\beta+\gamma}\circ
 Rq_{\alpha,\beta+\gamma*}p_{\alpha,\beta+\gamma}^*
 (1\boxtimes c_{\beta,\gamma})\circ C_R.
 \end{aligned}
\]

\begin{theorem}[Associativity from coefficient compatibility]
\label{thm:hall-coefficient-compatibility}
Under the preceding category, properness and base-change hypotheses, assume
$L=R$ as derived morphisms for every triple of charges. Then $\star$ is
associative on $\bigoplus_\gamma\mathbb H^*(\M_\gamma,K_\gamma)$.
For one fixed triple, the equality gives associativity on those three input
components alone.
\end{theorem}
\begin{proof}
The external product of three classes pulls back along $s$ to a class in
$\mathbb H^*(\mathcal F_{\alpha,\beta,\gamma},
s^*(K_\alpha\boxtimes K_\beta\boxtimes K_\gamma))$.
By derived global sections under $r$, this is a class in $\mathbb H^*(D)$.
Naturality of pullback, the external product and the two base-change
identifications sends this class to the two successive pullbacks defining
$(a\star b)\star d$ and $a\star(b\star d)$. Applying the coefficient maps
therefore gives $\mathbb H(L)$ and $\mathbb H(R)$ on the same class. Their
equality proves the assertion. No factors have been interchanged, so no
Koszul sign occurs in this associativity comparison. Applying the argument
to each triple and using bilinearity proves the direct-sum statement.
\end{proof}

For the matrix potential $\operatorname{Tr}(A[B,C])$, the coefficient maps
and identities of Chapters~\ref{ch:two-nonzero-block-potentials}
and~\ref{fl221:ch} are constructed on the full critical loci using
vanishing cycles. Their block order is submodule first. The precise order
conversion and the triples covered by these constructions are given in
Chapter~\ref{ch:extension-correspondences}.
